\documentclass[12pt,a4paper]{article}

\usepackage[top=1in, bottom=1in, left=1in, right=1in]{geometry}
\usepackage{amsmath}
\usepackage{booktabs}
\usepackage{multirow}
\usepackage{parskip}
\usepackage{titlesec}
\usepackage{hyperref}
\usepackage{array}
\usepackage{xcolor}
\usepackage{caption}

\titleformat{\section}{\large\bfseries}{}{0em}{}[\titlerule]
\titleformat{\subsection}{\normalsize\bfseries}{}{0em}{}

\hypersetup{colorlinks=true, linkcolor=blue, urlcolor=blue}

\begin{document}

% ── Title ─────────────────────────────────────────────────────────────────────
\begin{center}
    {\Large\bfseries University of Bahrain}\\[4pt]
    College of Information Technology\\
    Department of Computer Science\\[6pt]
    {\large ITCS 440: Intelligent Systems}\\
    Second Semester 2025/2026\\[8pt]
    \rule{\linewidth}{0.8pt}\\[6pt]
    {\LARGE\bfseries Assignment 2: Neural Networks}\\[4pt]
    {\large Backpropagation Experiments on UCI Datasets}\\[6pt]
    \rule{\linewidth}{0.8pt}\\[10pt]
\end{center}

\noindent
\begin{tabular}{|l|p{9cm}|}
\hline
\multirow{2}{*}{\textbf{Student Names}} & Abdulnaser Fawzi Hamadi \\
\cline{2-2}
 & Khalid Jassas Saleh \\
\hline
\multirow{2}{*}{\textbf{Student IDs}} & 202008866 \\
\cline{2-2}
 & 202008213 \\
\hline
\textbf{Section}      & 3 \\
\hline
\textbf{Submission Date} & May 17, 2026 \\
\hline
\end{tabular}

\bigskip

% ═════════════════════════════════════════════════════════════════════════════
\section{Introduction}

In this assignment I used a backpropagation neural network to classify two
different datasets taken from the UCI Machine Learning Repository. The idea
behind backpropagation is straightforward: the network makes a prediction,
measures how wrong it was, then works backwards through each layer adjusting
the weights slightly to reduce that error. This process repeats for many
iterations until the model stops improving.

Three parameters were changed between runs to see how they affect the outcome.
The \textbf{learning rate} decides how big each weight update is. Too large and
the model jumps past the solution; too small and it barely moves.
\textbf{Momentum} carries some of the previous update forward into the next
step, which helps when the gradient keeps pointing in the same direction.
\textbf{Epochs} is simply how many times the algorithm goes over the full
training data. More passes usually mean better accuracy, up to a point.

I ran four different combinations of these parameters on each dataset and
recorded the accuracy and F1-score for every run. The goal was to observe
how changing one parameter at a time shifts the results and to find out which
combination works best.

% ═════════════════════════════════════════════════════════════════════════════
\section{Datasets}

\subsection{Dataset 1: Pima Indians Diabetes (UCI)}

\textbf{Source:} UCI Machine Learning Repository\\
\textbf{URL:} \url{https://archive.ics.uci.edu/ml/datasets/diabetes}\\
\textbf{Instances:} 768\\
\textbf{Features:} 8 numeric attributes (pregnancies, glucose, blood pressure,
skin thickness, insulin, BMI, diabetes pedigree function, age)\\
\textbf{Classes:} 2 (positive/negative diabetes diagnosis)\\
\textbf{Task:} Binary classification

This dataset contains medical records from female patients and the task is to
predict whether a patient has diabetes or not. I picked it because it is
well-studied, small enough to qualify under 1000 instances, and the features
are straightforward medical measurements. One thing worth noting is that some
values like insulin and glucose are recorded as zero in the file, which
actually represents missing data. This makes the classification harder than
it looks on paper.

\subsection{Dataset 2: Adult Census Income (UCI)}

\textbf{Source:} UCI Machine Learning Repository\\
\textbf{URL:} \url{https://archive.ics.uci.edu/ml/datasets/adult}\\
\textbf{Instances:} 48,842\\
\textbf{Features:} 14 attributes (age, workclass, education, marital status,
occupation, relationship, race, sex, capital gain/loss, hours per week, country)\\
\textbf{Classes:} 2 (income $\leq$50K or $>$50K)\\
\textbf{Task:} Binary classification

This dataset comes from the 1994 US Census and the goal is to predict whether
a person earns more or less than \$50,000 per year. It has over 48,000 records
which makes it a good choice for the large dataset requirement. The features
are a mix of numbers and categories like occupation and education level, so
the categorical ones had to be converted to numbers before the network could
use them. Despite having more features than the Pima dataset, the class
boundary here seems easier for the network to learn.

% ═════════════════════════════════════════════════════════════════════════════
\section{Experimental Setup}

For both datasets I used a single hidden layer with 100 neurons and the
sigmoid activation function. The network was trained using stochastic gradient
descent, which is the same algorithm used inside the WEKA MultilayerPerceptron
classifier. Each dataset was split 80/20 for training and testing and the
features were scaled before training so that no single feature dominates due
to its range.

Four runs were done per dataset. Each run kept the architecture the same and
only changed the momentum, learning rate, and number of epochs.

% ═════════════════════════════════════════════════════════════════════════════
\section{Results}

\subsection{Dataset 1: Pima Indians Diabetes}

\begin{table}[h]
\centering
\caption{Backpropagation results on Pima Indians Diabetes dataset (768 instances)}
\label{tab:diabetes}
\begin{tabular}{ccccc}
\toprule
\textbf{Momentum} & \textbf{Learning Rate} & \textbf{Epochs} & \textbf{Accuracy} & \textbf{F1-Score} \\
\midrule
0.2 & 0.10 & 100 & 68.83\% & 69.47\% \\
0.5 & 0.30 & 200 & 77.92\% & 76.03\% \\
0.9 & 0.01 & 500 & 74.03\% & 74.21\% \\
0.2 & 0.50 & 300 & 72.73\% & 73.14\% \\
\bottomrule
\end{tabular}
\end{table}

\subsection{Dataset 2: Adult Census Income}

\begin{table}[h]
\centering
\caption{Backpropagation results on Adult Census Income dataset (48,842 instances)}
\label{tab:adult}
\begin{tabular}{ccccc}
\toprule
\textbf{Momentum} & \textbf{Learning Rate} & \textbf{Epochs} & \textbf{Accuracy} & \textbf{F1-Score} \\
\midrule
0.2 & 0.10 & 100 & 85.46\% & 84.96\% \\
0.5 & 0.30 & 200 & 85.76\% & 85.31\% \\
0.9 & 0.01 & 500 & 85.62\% & 85.22\% \\
0.2 & 0.50 & 300 & 85.33\% & 84.18\% \\
\bottomrule
\end{tabular}
\end{table}

% ═════════════════════════════════════════════════════════════════════════════
\section{Analysis and Discussion}

\subsection{Dataset 1: Pima Indians Diabetes}

The Pima dataset results varied quite a bit between runs, which was expected
given how small the dataset is. Small datasets are more sensitive to how the
network is configured.

\textbf{Run 1} gave the worst result at 68.83\%. Only 100 epochs with a low
momentum of 0.2 meant the network stopped learning too early. There simply
wasn't enough training time, and without strong momentum the updates were too
weak to push the weights toward a good solution.

\textbf{Run 2} was clearly the best: 77.92\% accuracy and 76.03\% F1-score.
The learning rate of 0.3 is high enough that the weights update meaningfully
each step, and a momentum of 0.5 keeps things moving without getting unstable.
Going to 200 epochs gave it enough time to actually converge. In my opinion
this run worked because none of the three values were at an extreme.

\textbf{Run 3} had 500 epochs but a very low learning rate of 0.01, which
explains the 74.03\% result. The network was learning, just extremely slowly.
The high momentum of 0.9 helped carry it forward but even then 500 passes
weren't quite enough to close the gap. It's a case where increasing epochs
even further might have helped.

\textbf{Run 4} had the highest learning rate at 0.5 and got 72.73\%. A rate
that high tends to overshoot and the model ended up bouncing around instead of
settling. Adding more epochs (300) helped a little but the damage from early
unstable updates was hard to undo.

The main takeaway from this dataset is that extremes don't work well here.
The best run was the one where all three parameters were at moderate values.

\subsection{Dataset 2: Adult Census Income}

The Adult dataset behaved very differently from Pima. All four runs landed
between 85.33\% and 85.76\%, which is a much narrower range. When you have
nearly 40,000 training samples, the network gets enough information each epoch
that minor differences in learning rate or momentum don't matter as much.

\textbf{Run 1} reached 85.46\% accuracy even with just 100 epochs. The large
amount of data compensates for the fewer training passes, so the model still
learns a solid decision boundary.

\textbf{Run 2} was again the top performer at 85.76\% accuracy and 85.31\%
F1-score. The same mid-range settings that worked on the small dataset worked
here too.

\textbf{Run 3} got 85.62\%. The low learning rate slowed things down but the
combination of high momentum and 500 epochs recovered the loss. It essentially
took longer to get to a similar place.

\textbf{Run 4} is interesting: the accuracy was fine at 85.33\% but the
F1-score dropped noticeably to 84.18\%. The high learning rate (0.5) made the
model less precise on the minority class, which in this dataset is people
earning over \$50K. Accuracy alone hides this problem because the majority
class is large enough to keep the percentage high regardless.

\subsection{Comparison Between Datasets}

The two datasets are quite different in how they respond to parameter changes.
Pima showed a 9-point swing in accuracy across runs while Adult barely moved 0.5
points. The reasons come down to three things. First, more data means more stable
training. Second, the Adult features like education and occupation carry stronger
signals for the prediction task than the Pima medical measurements which have
noisy zeros. Third, the Pima minority class (diabetic patients) is harder to
separate from the majority class given only 8 features.

One consistent pattern is that Run 2 came out on top for both datasets.
This suggests that mid-range values tend to be safer than extremes regardless
of how big the dataset is.

% ═════════════════════════════════════════════════════════════════════════════
\section{Conclusion}

Running these experiments on two datasets of very different sizes gave a clear
picture of how parameter tuning affects backpropagation. On the Pima dataset the
difference between a bad run and the best run was almost 10 percentage points,
which shows how much the settings matter when training data is limited. On the
Adult dataset the results were much more stable and the network performed well
across all four configurations.

The setting that worked best in both cases was a learning rate of 0.3, momentum
of 0.5, and 200 epochs. Nothing extreme. In practice this makes sense: a
moderate learning rate takes steps that are large enough to make real progress
but small enough not to overshoot, and momentum of 0.5 smooths out the updates
without over-relying on past gradients.

From these runs I'd say the most important lesson is to watch both accuracy and
F1-score together. The high learning rate run on the Adult dataset looked fine
by accuracy but the F1 told a different story. For problems with class imbalance,
F1 is the more honest measure of how well the model actually works.

\end{document}
